% Generated draft: requires mathematical and source review.
\documentclass[11pt]{article}
\usepackage[margin=25mm]{geometry}
\usepackage{fontspec}
\setmainfont{Latin Modern Roman}
\usepackage{amsmath,amssymb,mathtools,unicode-math}
\setmathfont{Latin Modern Math}
\usepackage{enumitem,needspace,xurl,hyperref}
\setlength{\parindent}{0pt}
\setlength{\parskip}{5pt}
\setlength{\emergencystretch}{3em}
\newcommand{\N}{\mathbb{N}}
\newcommand{\Z}{\mathbb{Z}}
\newcommand{\Q}{\mathbb{Q}}
\newcommand{\R}{\mathbb{R}}
\newcommand{\C}{\mathbb{C}}
\newcommand{\F}{\mathbb{F}}
\newcommand{\A}{\mathbb{A}}
\newcommand{\PP}{\mathbb{P}}
\newcommand{\E}{\mathbb{E}}
\newcommand{\eps}{\varepsilon}
\newcommand{\dd}{\,\mathrm{d}}
\begin{document}
\section*{UnsolvedMath ID: 1900002}
\textbf{Source problem code:} catalogue entry
\subsection{Definitions and mathematical statement}
An integer polygon is a polygon in \(\R^2\) with vertices in \(\Z^2\). For distinct lattice points \(A,B\), its lattice edge length is
\[
\ell(AB)=\gcd(|B_1-A_1|,|B_2-A_2|).
\]
For an angle between primitive integer vectors \(u,v\), its lattice sine is \(|\det(u,v)|\). Under the usual sail convention for a pointed rational convex angle with lattice vertex \(O\), take the finite polygonal boundary facing \(O\) of the convex hull of all nonzero lattice points in its cone, translated so \(O=0\). If its successive vertices are \(B_0,\ldots,B_m\), the lattice-length--sine sequence (LLS-sequence) alternates its edge lengths and interior lattice sines:
\[
\bigl(\ell(B_0B_1),\operatorname{lsin}\angle B_0B_1B_2,
\ell(B_1B_2),\ldots,\ell(B_{m-1}B_m)\bigr).
\]
Here the determinant for an interior lattice sine uses the primitive vectors along the two incident edges. Endpoint order depends on the orientation convention.

The IKEA problem asks, for each polygon size \(n\ge3\), to characterize exactly the tuples \((\sigma_1,\ldots,\sigma_n)\) that occur as the LLS-sequences of the successive angles of an integer \(n\)-gon. In a convex-polygon version using the preceding convention, realizability means that there exist cyclically ordered lattice vertices \(v_1,\ldots,v_n\) of a nondegenerate convex polygon with exactly \(n\) corners such that
\[
\operatorname{LLS}(\angle v_{i-1}v_iv_{i+1})=\sigma_i
\qquad(1\le i\le n),
\]
with indices modulo \(n\). The requested classification would give necessary and sufficient compatibility conditions on the entire tuple, not merely characterize individual lattice angles.

The excerpt does not provide its LLS convention or state whether polygons must be convex, merely simple, or may self-intersect, nor which cyclic/reversal identifications of tuples are allowed. The convex version above is therefore explicit conditional notation, not a silent restriction of the source's full target. The supplied research summary labels this OPEN-TRIAGE with moderate confidence, mentions unspecified constructive partial results, and says a conclusive realization/classification result could not be verified.
\subsection{Short English statement}
Determine exactly which collections of lattice angle LLS-sequences fit together as the successive angles of a lattice polygon. The polygon class and sequence conventions require confirmation from the source.
\subsection{Sources}
\begin{itemize}
\item Local UnsolvedMath registry, permanent ID \(1900002\), exact source code \texttt{catalogue entry}, ``Geometry of Continued Fractions --- Integer trigonometry and IKEA problem.''
\item Karpenkov, \emph{Open problems in geometry of continued fractions} (2017), Problem 2, ``IKEA problem''; \url{https://arxiv.org/abs/1712.01450}. Extraction from source TeX; recorded access 2026-07-29.
\item Supplied moderate-confidence OPEN-TRIAGE summary, without specific references for its mentioned constructive partial results.
\end{itemize}
\end{document}
